\section{Nouvelle approche de r\'esolution dans l'espace des moments}

\begin{frame}
 \frametitle{Progression}
 \framesubtitle{}
 \begin{itemize}
   \item \textcolor{gray}{Introduction}
   \item \textcolor{gray}{G\'en\'eralit\'es sur l'interaction laser-atome 
en champ fort}
   \item \textbf{Nouvelle approche de r\'esolution 
dans l'espace des moments}
   \item R\'esultats
   \item Conclusion
 \end{itemize}
\end{frame}



\begin{frame}
 \frametitle{Nouvelle approche de r\'esolution dans l'espace des moments}
 \framesubtitle{Fonctions sturmiennes dans l'\acrshort{ec}}

%\begin{tikzpicture}
% \node at (0,0) {};
% \node at (12,5) {\includegraphics[width=.45\linewidth]{../figures/chap3/sturm_conf_variation_kappa.eps}};
%% \node at (7,.5) {\includegraphics[width=.7\linewidth]{../figures/chap3/sturm_mom_variation_nn_nbar.eps}};
%\end{tikzpicture}

\begin{textblock*}{10pt}(265pt,60pt)
\includegraphics[scale=.3]{../figures/chap3/sturm_conf_variation_kappa.eps}
\end{textblock*}

\begin{textblock*}{10pt}(30pt,60pt)
\footnotesize
\begin{equation*}
 \begin{split} 
  & \Phi_{n\ell m}^\kappa (\boldsymbol{r}) = \frac{1}{r} S_{n\ell}^\kappa (r)
   Y_{\ell m}(\theta,\varphi), \qquad
    \ell=0,1,2,3,... \qquad n\geqslant \ell+1, \\
   &  r\Phi_{n\ell m}^\kappa (\boldsymbol{r}) \xrightarrow{r \rightarrow 0} 0 \qquad ; \qquad r\Phi_{n\ell m}^\kappa (\boldsymbol{r})  \xrightarrow{r \rightarrow \infty} 0 \qquad
 \end{split}
\end{equation*}

\end{textblock*}


\begin{textblock*}{10pt}(30pt,100pt)
\footnotesize
\begin{equation*}
 \begin{split} 
&  S_{n\ell}^\kappa (r) = \widetilde{N}_{n\ell}^\kappa \, r^{\ell+1}
    {\rm e}^{-\kappa r} {}_1 F_1(n-\ell-1;2\ell+1;2\kappa r), \\
& \widetilde{N}_{n\ell}^\kappa = \sqrt{\frac{\kappa}{n}}(2\kappa)^{\ell+1}
    \frac{1}{(2\ell+1)!}\sqrt{\frac{(n+\ell)!}{(n-\ell-1)!}}.
 \end{split}
\end{equation*}

\end{textblock*}



\begin{textblock*}{100pt}(30pt,180pt)
\footnotesize
Relation de r\'ecurrence
\begin{equation*}
 \begin{split}
 \sqrt{n(n-\ell-1)(n+\ell)}S_{n,\ell}^\kappa(r)
 - 2 &\sqrt{(n-1)}(n-1-\kappa r) S_{n-1,\ell}^\kappa(r) \\
 - & \sqrt{(n-2)(n-\ell-2)(n+\ell+1)}S_{n-2,\ell}^\kappa(r) = 0.
 \end{split}
\end{equation*}

\end{textblock*}


% \begin{itemize}
%   \item Fonctions sturmiennes dans l'\acrshort{ec} et lien avec 
%les fonctions propres des hydrog\'eno\"ides (courbes)
%    \item Int\'er\^et math\'ematique des fonctions sturmiennes  
%    \item Fonctions sturmiennes dans l'\acrshort{em} (courbes)
% \end{itemize}

\end{frame}


\begin{frame}
 \frametitle{Nouvelle approche de r\'esolution dans l'espace des moments}
 \framesubtitle{Fonctions sturmiennes dans l'\acrshort{em}}

%\begin{tikzpicture}
% \node at (0,0) {};
% \node at (12,5) {\includegraphics[width=.45\linewidth]{../figures/chap3/sturm_mom_variation_kappa.eps}};
%% \node at (7,.5) {\includegraphics[width=.7\linewidth]{../figures/chap3/sturm_mom_variation_nn_nbar.eps}};
%\end{tikzpicture}

\begin{textblock*}{10pt}(265pt,60pt)
\includegraphics[scale=.3]{../figures/chap3/sturm_mom_variation_kappa.eps}
\end{textblock*}

\begin{textblock*}{10pt}(30pt,60pt)
\footnotesize
\begin{equation*}
 \begin{split} 
  & \Phi_{n\ell m}^\kappa (\boldsymbol{p})
      = \Sigma_{n\ell}^\kappa (p)
      \widetilde{Y}_{\ell m}(\hat{\boldsymbol{p}}),  \qquad
\widetilde{Y}_{\ell m}(\hat{\boldsymbol{p}})
= (-1)^\ell Y_{\ell m}(\hat{\boldsymbol{p}}), \\
   & \Sigma_{n\ell }^\kappa (p)
   = \sqrt{\frac{2}{\pi}} \int dr\, rj_\ell(pr) S_{n\ell}^\kappa (r).
 \end{split}
\end{equation*}

\end{textblock*}


\begin{textblock*}{10pt}(30pt,100pt)
\footnotesize
\begin{equation*}
 \begin{split}
 & \Sigma_{n\ell}^\kappa (p)
    = \mathcal{N}_{n\ell}^\kappa \frac{p^\ell}{(\kappa^2 + p^2)^{\ell+2}}
      C_{n-\ell-1}^{(\ell+1)}
  \left(\frac{p^2 - \kappa^2}{p^2 + \kappa^2} \right), 
 \;\;\;\; t=\left( \frac{p^2 - \kappa^2}{p^2 + \kappa^2} \right), \\
  & \mathcal{N}_{n\ell}^\kappa = 2^{2\ell+2}\kappa^{\ell+2}\,\ell!\,
      \sqrt{\frac{2\kappa n}{\pi}\frac{(n-\ell-1)!}{(n+\ell)!}}.
 \end{split}
% \begin{split} 
%& \Sigma_{n\ell}^\kappa (t)
%    = \widetilde{\mathcal{N}}_{n\ell}^\kappa\,(1-t)^{\ell/2+2} (1+t)^{\ell/2}\,
%      C_{n-\ell-1}^\kappa \left( t \right),
% \;\;\;\;\; t=\left( \frac{p^2 - \kappa^2}{p^2 + \kappa^2} \right), \\
%  & \widetilde{\mathcal{N}}_{n\ell}^\kappa = \frac{2^\ell\,\ell!}{\kappa^2}\,
%      \sqrt{\frac{2\kappa n}{\pi}\frac{(n-\ell-1)!}{(n+\ell)!}}.
% \end{split}
\end{equation*}

\end{textblock*}


\begin{textblock*}{100pt}(30pt,180pt)
\footnotesize
Relation de r\'ecurrence
\begin{equation*}
 \begin{split}
\sqrt{(n-2)(n-\ell-1)(n+\ell)} \, & \Sigma_{n,\ell}^\kappa (p)
   + \sqrt{n(n-1)(n-2)} \left( \frac{p^2 -\kappa^2}{p^2+\kappa^2} \right)
     \Sigma_{n-1,\ell}^\kappa (p) \\
    & \hspace{25mm} + \sqrt{n(n-\ell-2)(n+\ell-1)} \Sigma_{n-2,\ell}^\kappa (p) = 0.
 \end{split}
\end{equation*}

\end{textblock*}


% \begin{itemize}
%   \item Fonctions sturmiennes dans l'\acrshort{ec} et lien avec 
%les fonctions propres des hydrog\'eno\"ides (courbes)
%    \item Int\'er\^et math\'ematique des fonctions sturmiennes  
%    \item Fonctions sturmiennes dans l'\acrshort{em} (courbes)
% \end{itemize}

\end{frame}


\begin{frame}
 \frametitle{Nouvelle approche de r\'esolution dans l'espace des moments}
 \framesubtitle{Potentiel coulombien non local et singularit\'e logarithmique}
 
% \putat{-0}{-120}{\includegraphics[scale=0.25]{../figures/chap3/sturm_conf_variation_nn_nbar.eps}}
% \putat{-0}{0}{Equation de Schr\"odinger dans l'\acrshort{ec} : 
%  \begin{equation}
%    i\frac{\partial}{\partial t}\Psi(\boldsymbol{r},t) =
% \end{equation}
%}


% \begin{itemize}
%   \item Equation de Schr\"odinger dans l'\acrshort{ec}
%   \begin{equation}
%     i\frac{\partial}{\partial t}S_{n\ell}(r,t) 
%   = \left(\frac{\boldsymbol{p}^2}{2} -\frac{Z}{r} + \boldsymbol{\mathcal{E}}\cdot\boldsymbol{r}\right)S_{n\ell}(r,t)
%  \end{equation}
%   \item Obtention de l'\'equation de Schr\"odinger dans l'\acrshort{em}
%   \item Probl\`eme du potentiel coulombien non local
% \end{itemize}


%\begin{textblock*}{10pt}(0pt,0pt)
%Hello
%\end{textblock*}


\begin{textblock*}{200pt}(30pt,60pt)
\footnotesize
\begin{itemize}
\item Equation de Schr\"odinger dans l'\acrshort{ec}
\end{itemize}
\begin{equation*}
 \left(\frac{d^2}{dr^2} - \frac{\kappa^2}{2} -\frac{\boldsymbol{L}^2}{2r^2}
   - \frac{\kappa n}{r}\right) \varphi (\boldsymbol{r}) = 0, 
\qquad \varphi(\boldsymbol{r}) = \sum_{n\ell m} c_{n\ell m}^\kappa \frac{1}{r}
 S_{n\ell}^\kappa (r) Y_{\ell m}({\boldsymbol{\hat{r}}}),
\qquad \kappa=\frac{Z}{n},\label{eqn:ss_config}
\end{equation*}

\end{textblock*}

\begin{textblock*}{200pt}(30pt,100pt)
\footnotesize
\begin{itemize}
\item Equation de Schr\"odinger dans l'\acrshort{em}
\end{itemize}
\begin{equation*}
 \begin{split}
& \left( \frac{\boldsymbol{p}^2}{2} - E \right) \psi(\boldsymbol{p})
   + \int d\boldsymbol{p}^\prime\,V(\boldsymbol{p},\boldsymbol{p}^\prime)\,
     \psi(\boldsymbol{p^\prime}) =  0, \qquad
\qquad \psi(\boldsymbol{p}) = \sum_{n\ell m} c_{n\ell m}^\kappa 
 \Sigma_{n\ell}^\kappa (p) \widetilde{Y}_{\ell m}({\boldsymbol{\hat{p}}}), \\
&\qquad\qquad  V(\boldsymbol{p},\boldsymbol{p^\prime}) 
  = -\frac{Z}{2\pi^2\left| \boldsymbol{p^\prime}-\boldsymbol{p}\right|^2}.
      \label{eqn:schrodinger_moment}
\end{split}
\end{equation*}

\end{textblock*}

\begin{textblock*}{100pt}(30pt,165pt)
\footnotesize
On montre\footnote{\Tiny 
Ongonwou \textit{et al.}, Ann. Phys. 375, 471–490, (2016)} que :
\end{textblock*}

\begin{textblock*}{100pt}(30pt,170pt)
\footnotesize
\begin{equation*}
   \frac{1}{\left| \boldsymbol{p^\prime}-\boldsymbol{p}\right|^2}
  = \sum_{\ell=0}^\infty \sum_{m=-\ell}^\ell \sum_{n=\ell+1}^\infty
   \frac{2\pi^2\kappa}{n} \widetilde{\Phi}_{n\ell m}^\kappa (\boldsymbol{p}) 
   \widetilde{\Phi}_{n\ell m}^{\kappa\ast} (\boldsymbol{p^\prime}),
 \qquad\qquad \widetilde{\Phi}_{n\ell m}^\kappa (\boldsymbol{p})= \frac{\kappa^2+p^2}{2\kappa^2}{\Phi}_{n\ell m}^\kappa (\boldsymbol{p}),
\end{equation*}
\end{textblock*}


\begin{textblock*}{100pt}(190pt,195pt)
\small
\begin{equation*}
   \left(\frac{\boldsymbol{p}^2}{2}-E\right) \psi(\boldsymbol{p})
  - \sum_{n\ell m} \frac{Z\kappa}{n} \widetilde{\Phi}_{n\ell m}^\kappa (\boldsymbol{p}) \int d\boldsymbol{p^\prime} \widetilde{\Phi}_{n\ell m}^{\kappa\ast} (\boldsymbol{p^\prime}) \psi(\boldsymbol{p^\prime}) = 0.
\end{equation*}
\end{textblock*}

\end{frame}


\begin{frame}
 \frametitle{Nouvelle approche de r\'esolution dans l'espace des moments}
 \framesubtitle{Formulation matricielle de l'\acrshort{esit} 
sous le nouveau mod\`ele dans la base sturmienne}

\begin{textblock*}{10pt}(30pt,60pt)
\footnotesize
\begin{equation*}
  \mathbf{H}^{(0)} \mathbf{\Psi} = E \mathbf{S}\mathbf{\Psi},
\end{equation*}
\end{textblock*}


\begin{textblock*}{200pt}(230pt,65pt)
\small
\begin{itemize}
 \item Probl\`eme aux valeurs propres g\'en\'eralis\'e
\end{itemize}
\end{textblock*}

\begin{textblock*}{10pt}(30pt,80pt)
\footnotesize
\begin{align*}
S_{(n^\prime \ell^\prime m^\prime),(n\ell m)}
  & = \left\langle n^\prime \ell^\prime m^\prime ; \kappa
   | n\ell m ; \kappa \right\rangle = Q_{n^\prime n}^{\kappa,\ell}
   \delta_{\ell^\prime \ell}\delta_{m^\prime m},  \\
 H^{(0)}_{(n^\prime \ell^\prime m^\prime),(n\ell m)}
  & = \mathcal{T}_{(n^\prime \ell^\prime m^\prime),(n\ell m)}
    + \mathcal{V}_{(n^\prime \ell^\prime m^\prime),(n\ell m)},
     \label{eqn:unperturbed_terms}
\end{align*}

\end{textblock*}


\begin{textblock*}{10pt}(30pt,130pt)
\footnotesize
\begin{align*}
& Q_{n^\prime,n}^{\kappa,\ell} = \delta_{n^\prime,n}
- \frac{1}{2}\sqrt{\frac{(n+\ell)(n-\ell-1)}{n(n-1)}} \delta_{n^\prime, n-1}
- \frac{1}{2}\sqrt{\frac{(n+\ell+1)(n-\ell)}{n(n+1)}} \delta_{n^\prime, n+1}. \\
&  \mathcal{T}_{(n^\prime \ell^\prime m^\prime),(n\ell m)}
%     = \Big\langle n^\prime \ell^\prime m^\prime ; \kappa
%  |\frac{p^2}{2} | n\ell m ; \kappa \Big\rangle, \hspace{8cm} \\
    \hspace{0cm} = \frac{\kappa^2}{2}\Bigg(\delta_{n^\prime n} 
    + \frac{1}{2}\sqrt{\frac{(n+\ell)(n-\ell-1)}{n(n-1)}}\delta_{n^\prime, n-1}
   \hspace{0cm} + \frac{1}{2}\sqrt{\frac{(n+\ell+1)(n-\ell)}{n(n+1)}}
   \delta_{n^\prime, n+1}\Bigg)\delta_{\ell^\prime \ell}\delta_{m^\prime m},\\
&  \mathcal{V}_{(n^\prime \ell^\prime m^\prime),(n\ell m)}
   = -\frac{Z\kappa}{n}\delta_{n^\prime n}
      \delta_{\ell^\prime \ell}\delta_{m^\prime m}.
\end{align*}

%\begin{equation}
% \begin{split}
%&  \mathcal{T}_{(n^\prime \ell^\prime m^\prime),(n\ell m)}
%      = \Big\langle n^\prime \ell^\prime m^\prime ; \kappa
%   |\frac{p^2}{2} | n\ell m ; \kappa \Big\rangle, \hspace{8cm} \\
%   & \hspace{2cm}  = \frac{\kappa^2}{2}\Big(\delta_{^\prime n} +
%   + \frac{1}{2}\sqrt{\frac{(n+\ell)(n-\ell-1)}{n(n-1)}}
%       \delta_{n^\prime, n-1}  \\
%    & \hspace{4cm} + \frac{1}{2}\sqrt{\frac{(n+\ell+1)(n-\ell)}{n(n+1)}}
%         \delta_{n^\prime, n+1}\Big)\delta_{\ell^\prime \ell}\delta_{m^\prime m},
%    \label{eqn:kinetic_sturmian}
%\end{split}
%\end{equation

\end{textblock*}

% \begin{itemize}
%   \item Equation de Schr\"odinger
%   \item Expressions des matrices de l'\'energie cin\'etique, 
%de l'\'energie potentielle et de recouvrement 
% \end{itemize}
\end{frame}

\begin{frame}
 \frametitle{Nouvelle approche de r\'esolution dans l'espace des moments}
 \framesubtitle{Formulation matricielle de l'\acrshort{esdt} 
sous le nouveau mod\`ele dans la base sturmienne}

\begin{textblock*}{10pt}(30pt,60pt)
\footnotesize
\begin{equation*}
\mathbf{H} = \mathbf{H}^{(0)} + \mathbf{H}^{(1)}
\qquad\qquad  H^{(1)} (t) = {A}(t)\cdot p_z,
\end{equation*}
\end{textblock*}

\begin{textblock*}{10pt}(40pt,80pt)
\footnotesize
\begin{equation*}
 \begin{split}
   H^{(1)}_{(n^\prime\ell^\prime m^\prime),(n\ell m)} 
  = A(t) \Bigg\lbrace
    \sqrt{\frac{(\ell+m+1)(\ell-m+1)}{(2\ell+1)(2\ell+3)}}
  \delta_{\ell^\prime,\ell+1}  \mathcal{A}_{\bar n^\prime,\bar n} 
  - \sqrt{\frac{(\ell+m)(\ell-m)}{(2\ell-1)(2\ell+1)}}
   \delta_{\ell^\prime,\ell-1} \mathcal{B}_{\bar n^\prime,\bar n}
  \Bigg\rbrace \delta_{m^\prime,m},
 \end{split}
\end{equation*}
\end{textblock*}

\begin{textblock*}{10pt}(40pt,120pt)
\footnotesize
\begin{align*}
& \mathcal{A}_{\bar n^\prime, \bar n}
 = \frac{\kappa}{2} \Bigg\lbrace
 \sqrt{\frac{(\bar n +2\ell+2)(\bar n+2\ell+3)}{(\bar n +\ell+1)(\bar n+\ell+2)}}\delta_{\bar n^\prime,\bar n}
 - \sqrt{\frac{\bar n(\bar n-1)}{(\bar n +\ell)(\bar n+\ell+1)}}\delta_{\bar n^\prime,\bar n-2} \Bigg\rbrace, \qquad\qquad \bar{n}=n-\ell-1, \\  
& \mathcal{B}_{\bar n^\prime, \bar n}
  = \frac{\kappa}{2} \Bigg\lbrace
 \sqrt{\frac{(\bar n +2\ell)(\bar n+2\ell+1)}{(\bar n +\ell)(\bar n+\ell+1)}}\delta_{\bar n^\prime,\bar n}
 - \sqrt{\frac{(\bar n+1)(\bar n+2)}{(\bar n +\ell+1)(\bar n+\ell+2)}}\delta_{\bar n^\prime,\bar n+2} \Bigg\rbrace. 
\end{align*}
\end{textblock*}


\begin{textblock*}{500pt}(40pt,200pt)
\footnotesize
\begin{itemize}
 \item Seuls les blocs $\Delta\ell=0,\pm 1$ sont occup\'es.
 \item Conservation du nombre quantique magn\'etique $m$ au cours de  
la propagation.
 \item \textbf{Manifestation des propri\'et\'es de l'approximation dipolaire}.
\end{itemize}
\end{textblock*}

\end{frame}

\begin{frame}
 \frametitle{Nouvelle approche de r\'esolution dans l'espace des moments}
 \framesubtitle{Sch\'ema de propagation Crank-Nicholson de Nurhuda 
et Faisal\footnote{\Tiny M. Nurhuda \textit{et al.}, Phys. Rev. A 60, 4, (1999)} (\acrshort{cnfa})}


\begin{textblock*}{100pt}(30pt,70pt)
\footnotesize
Dans la base sturmienne:
\begin{equation*}
i\mathbf{S}\frac{\partial}{\partial t}\mathbf{C} = \mathbf{H}\mathbf{C}, 
%\qquad\mathbf{H} = \mathbf{H}^{(0)} + \mathbf{H}^{(1)}
\end{equation*}

\end{textblock*}

\begin{textblock*}{100pt}(30pt,110pt)
\footnotesize
Dans la base atomique:
\begin{equation*} 
\footnotesize
\widetilde{\mathbf{H}} = \mathbf{U}^{-1} \mathbf{H}\;\mathbf{U}, 
  \qquad\widetilde{\mathbf{C}} =  \mathbf{U}^{-1}{\mathbf{C}} 
\end{equation*}
 
\end{textblock*}


\begin{textblock*}{200pt}(30pt,140pt)
\footnotesize
\begin{equation*}
 \widetilde{\mathbf{C}}(t+\Delta t)
   \sim \frac{\boldsymbol{\mathds{1}}-\xi\widetilde{\mathbf{H}}}
{\boldsymbol{\mathds{1}}+\xi\widetilde{\mathbf{H}}}\widetilde{\mathbf{C}}(t) \qquad\Longrightarrow\qquad  \widetilde{\mathbf{C}}(t+\Delta t)
\sim (\boldsymbol{\mathds{1}}+\xi\widetilde{\mathbf{H}})^{-1} \mathbf{\Phi},
\qquad \xi = i\frac{\Delta t}{2}
\end{equation*}

\end{textblock*}


\begin{textblock*}{10pt}(30pt,170pt)
\footnotesize
\begin{equation*}
\mathds{1}+\xi\widetilde{\mathbf{H}} = \mathbf{O_D} + \mathbf{O_{ND}},\qquad
\mathbf{O_D} = \mathbf{\mathds{1}}+\xi\widetilde{\mathbf{H}}^{(0)},
 \qquad \mathbf{O_{ND}} = \xi\widetilde{\mathbf{H}}^{(1)},
\end{equation*}

\end{textblock*}

\begin{textblock*}{10pt}(30pt,190pt)
\footnotesize
\begin{equation*}
 \widetilde{\mathbf{C}}(t+\Delta t) = 
 \Bigg\lbrace\Big[\mathds{1}-\mathbf{O_D^{-1}}\mathbf{O_{ND}} 
+\left(\mathbf{O_D^{-1}}\mathbf{O_{ND}}\right)^2
  +\ldots \Big]\mathbf{O_D}^{-1} \Bigg\rbrace \mathbf{\Phi}, 
  \qquad\qquad \big\| \widetilde{\mathbf{C}}(t+\Delta t) \big\| - 1 \leqslant \varepsilon
\end{equation*}

\end{textblock*}

%\begin{itemize}
%  \item Pr\'esentation du mod\`ele de propagation \acrshort{cnfa}
%  \item Impl\'ementation de la matrice d'interaction
%  \item Comparaison des performances Runge-Kutta - \acrshort{cnfa}

%\end{itemize}
\end{frame}

\begin{frame}
 \frametitle{Nouvelle approche de r\'esolution dans l'espace des moments}
 \framesubtitle{Calcul des observables}


\begin{textblock*}{200pt}(30pt,70pt)
\footnotesize
Paquet d'onde final dans la base sturmienne
\footnotesize
\begin{equation*}
   \Psi(\boldsymbol{p},\tau)
    = \sum_{n\ell} c_{n\ell}(\tau) \Sigma_{n\ell}(p)
       \widetilde{Y}_{\ell 0}(\hat{\boldsymbol{p}}),
\qquad\qquad \tau=n_{c}T.
    \label{eqn:final_wave}
\end{equation*}

\end{textblock*}


\begin{textblock*}{350pt}(30pt,105pt)
\begin{itemize}
\footnotesize
\item Populations des \'etats li\'es  :\hspace{30pt}
%$
% \Phi_i (\boldsymbol{p}) = \sum_{n\ell m} d_{n\ell}^{(i)}
%   \Phi_{n\ell m}^\kappa (\boldsymbol{p}).
%$
\begin{equation*}
\footnotesize
 \label{eqn:bound_state_pop}
 \left\langle \Phi_i(\boldsymbol{p}) | \Psi(\boldsymbol{p},t) \right\rangle
   = \sum_{\substack{n^\prime\ell^\prime m^\prime \\ n\ell m}}
    d_{n^\prime\ell^\prime}^{(i)\ast}  c_{n\ell m }(t)
  \left\langle \Phi_{n^\prime\ell^\prime m^\prime}^\kappa(\boldsymbol{p}) | \Phi_{n\ell m}^\kappa (\boldsymbol{p}) \right\rangle, 
 \qquad \mathcal{P}(t)=1-\sum_{n\ell m} d_{n\ell m}(t),
\end{equation*}
\end{itemize}
\end{textblock*}


\begin{textblock*}{10pt}(200pt,94pt)
\footnotesize
\begin{equation*}
  \Phi_i (\boldsymbol{p}) = \sum_{n\ell m} d_{n\ell}^{(i)}
    \Phi_{n\ell m}^\kappa (\boldsymbol{p}).
\end{equation*}
\end{textblock*}


\begin{textblock*}{400pt}(30pt,150pt)
\footnotesize
\begin{itemize}
\item Densit\'e de probabilt\'e :\hspace{50pt}
%\footnotesize
%$
%    \mathcal{D}(\boldsymbol{p},t)
%   = | \Psi(\boldsymbol{p},t) |^2
%    =  \left|  \sum_{n\ell} c_{n\ell}(t) \Sigma_{n\ell}(p)
%       \widetilde{Y}_{\ell 0}(\hat{\boldsymbol{p}})\right|^2.
%$
\end{itemize}
\end{textblock*}


\begin{textblock*}{500pt}(30pt,145pt)
\footnotesize\begin{equation*}
    \mathcal{D}(\boldsymbol{p},t)
   = | \Psi(\boldsymbol{p},t) |^2
    =  \left|  \sum_{n\ell} c_{n\ell}(t) \Sigma_{n\ell}(p)
       \widetilde{Y}_{\ell 0}(\hat{\boldsymbol{p}})\right|^2.
\end{equation*}
\end{textblock*}

\begin{textblock*}{200pt}(30pt,170pt)
\footnotesize
\begin{equation*}
 d^3\mathcal{P}
  = \mathcal{D}(\boldsymbol{p}) 2\pi p_n dp_n dp_z
  = \mathcal{D}(\boldsymbol{p}) pdEd\hat{\boldsymbol{p}},
\end{equation*}
\end{textblock*}


\begin{textblock*}{400pt}(30pt,190pt)
\begin{itemize}
\footnotesize
\item Distribution en moments : 
\item Distribution en \'energie :
\item Distribution en impulsions transversale et longitudinale.
\end{itemize}
\end{textblock*}


\begin{textblock*}{10pt}(300pt,170pt)
\footnotesize
\begin{align*}
   \frac{d\mathcal{P}}{dp_ndp_z}
&  = 2\pi p_n \mathcal{D}(\boldsymbol{p}), \\
   \frac{d\mathcal{P}}{dE}
&  =  \left|\sum_{n\ell m} c_{n\ell}^{\rm ion}\right|.
\end{align*}
\end{textblock*}


% \begin{itemize}
%   \item Evolution des populations au cours de l'interaction et 
%populations finales
%   \item Densit\'es d'\'energie
%   \item Spectres d'\'energie 
%   \item Distributions en moments et distributions angulaires
%   \item Distributions en impulsion 
% \end{itemize}
\end{frame}

